\section{Limitations and Future Work}
\label{sec:limitations}

\paragraph{Estimating safety-alignment strength.}
IAM-Teamer applies its soft-opening uniformly to every target because the attacker cannot reliably tell, in advance, how strongly a given model is aligned.
As the Llama analysis shows (Section~\ref{sec:analysis}), the same conservative opening that buys large gains on strongly aligned targets wastes an early turn on weakly aligned ones; an ideal attacker would instead calibrate its opening aggressiveness to the target's alignment strength.
Doing so requires a measure of that strength, and we were unable to identify a principled, model-inferable criterion that reliably distinguished strongly aligned targets from weakly aligned ones.
We considered the target's first-turn refusal rate (threshold $0.4$) as a proxy for alignment strength, but it is coarse, conflating robust alignment with merely stylistic refusal, shifting with the goal distribution, and offering no clear boundary separating ``strong'' from ``weak.''
Conditioning the opening on this proxy did not pay off.
It labels GPT-4o weak on the strength of a moderate first-turn refusal rate ($0.35$), but GPT-4o is only relatively weak, and switching it to the aggressive single-shot opening reserved for weak targets (Supplementary Section~\ref{app:attack_prompt}) leaves its ASR essentially unchanged from the uniform soft-opening ($98.0\%$ versus $99.5\%$ at $k{=}10$, and $44.5\%$ versus $51.5\%$ at $k{=}1$), since it still refuses about a third of overt first turns.
The proxy thus yields no gain even where it is most confident, while risking the turn-1 refusal cascade the soft-opening prevents wherever it misreads a resistant target as weak, so we discarded it and fall back to the uniform soft-opening.
A natural direction for future work is to estimate alignment strength online from richer interaction signals, such as early-turn refusal patterns, response consistency, or other behavioral fingerprints, so that the attacker can adaptively switch between a soft-opening and an overt first turn.

\paragraph{Dependence on strategy reuse.}
Warm-start is most effective when the attacker reuses and refines previously successful strategies.
However, because the ranked memory acts only as guidance rather than a constraint (Section~\ref{sec:strategy}), the attacker may instead generate new strategies from scratch.
In such cases, the seeded ranking provides little information to exploit, limiting the benefit of warm-start.
Future work could investigate mechanisms that adapt the strength of memory guidance or explicitly model when strategy invention is preferable to strategy reuse.

\paragraph{Dependence on the attack and judge models.}
IAM-Teamer's effectiveness is tied to the two LLMs that drive it, the attack model $A$ and the judge model $J$, both instantiated with Gemma-4-31B in our experiments.
The attacker bounds the diversity and quality of the strategies that can be applied, composed, or invented, so a weaker attacker explores a smaller and less effective strategy space.
The judge is even more central, since its score $r$ both decides attack success and supplies the binary reward for Thompson sampling. Any systematic bias or noise in the judge therefore propagates directly into strategy ranking and into the reported ASR; a judge that under- or over-credits particular response styles would mis-steer strategy selection and distort the comparison.
Because $A$ and $J$ are instantiated from the same model, the judge does more than score: it also coaches the attacker, distilling the success factors fed back to $A$. A blind spot shared between the two could thus inflate both the guidance and the reported ASR.
Although our judge validation (Section~\ref{sec:validation}) shows high agreement with an external classifier, this reliance on a single attacker--judge pair remains a limitation.
Future work could use a judge from a different model family than the attacker to break this loop, adopt ensemble judges to reduce scoring variance, and measure how sensitive IAM-Teamer's gains are to the choice of attacker and judge.